\chapter{Obstructions: what moving the problem cannot buy}\label{ch:obstructions}

The programme kept 148 leads against ECDLP in a lead atlas
(\art{crypto-autoresearcher:knowledge/leads/ecdlp/}, reviewed 2026-10-24).
Of these:
\begin{itemize}
\item 73 are closed with a named obstruction;
\item 51 are open with a decided next step;
\item 16 are blocked;
\item 8 are unexamined.
\end{itemize}
Most obstructions are textbook results applied to the case at hand: the
Shoup--Nechaev generic bound, the Corrigan-Gibbs--Kogan preprocessing bound
$ST^2\gtrsim N$~\cite{cgk2018}, isogeny invariance of the group order, and
the Balasubramanian--Koblitz embedding-degree bound. This chapter records
the few that the programme \emph{proved} itself and that constrain index
calculus specifically. Each states a lever that cannot work, and why.

\section{Transport conservation}\label{sec:obs-transport}

A homomorphic transfer $\varphi:E\to E'$ with rational kernel $\kappa$
changes the number of ordered $m$-decompositions of a target by exactly
$\abs{\kappa}^{m-1}$ (Lemma~\ref{lem:conservation} in
Chapter~\ref{sec:transfer}). So an isogeny buys only three things:
\begin{itemize}
\item the kernel symmetry, which is trivial on prime-order curves;
\item the factor-base geometry available on the destination curve;
\item the cost of arithmetic there.
\end{itemize}
Chains of isogenies compose to one homomorphism. A loop is an endomorphism
acting as a known scalar. Multi-curve factor bases glued by isogenies are
single-curve bases in disguise. Measured across sizes at fixed cofactor,
fibre-aware index calculus loses ground to rho as $\#E^{1/6}$; the full
measurements are in Chapter~\ref{sec:transfer}.

\section{The summation cover is the same everywhere}

Chapter~\ref{ch:cover} proved:
\begin{itemize}
\item the monodromy of the summation cover is $(\ZZ/2)^{m-2}$ for every
  curve;
\item its fibres split completely on the factor-base locus;
\item mean yield at fixed base size is a double count.
\end{itemize}
Together these close three families of leads: curve selection by
anomalous monodromy, fibre censuses as a relation-rate probe, and
factor-base geometries claimed to raise yield.

\section{Oracle-augmented index calculus stays generic}

Several leads proposed giving index calculus an extra oracle on the group:
\begin{itemize}
\item a jet or $p$-adic lift;
\item an endomorphism;
\item an incidence or halving test;
\item an elliptic-net evaluation.
\end{itemize}
The programme proved, at derivation tier, that each of these is
simulable in the generic group model, so it cannot beat
$\sqrt r$ (\texttt{KN-FIND-002}, \texttt{KN-FIND-b7e091},
\texttt{KN-FIND-194294}; halving is equivalent to an $x$-coordinate oracle).

The minimal non-simulable order-based identifier is a threshold oracle
$C_t$ (\texttt{KN-FIND-982fdf}). Table-based decomposition then obeys the
preprocessing bound $ST^2\gtrsim N$. Balanced generic index calculus costs
about $r^{2/3}/n^{4/3}$, consistent with the measured growth of the Koblitz
gap to rho as $r^{0.165\,[0.117,0.213]}$
(\texttt{LEAD-ICB-generic-table-oracle}). These are known instances of
generic-group reasoning~\cite{cgk2018}; they are recorded because each
retired a concrete proposal.

\section{Lifting sections does not produce relations}

Silverman's xedni calculus lifts points to a curve over a global field, in
the hope that rank forces relations. On a frozen family of rational
elliptic surfaces (with $\deg b=6$ and monic quadratic $x$), the programme
derived the exact lift probability
\[
  P_{\rm lift}(p)=\frac{1-2/p+p^{-3}}{2p^3},
\]
so the exponent is $\alpha=3$, far above the $\alpha<1/2$ that a relation
source would need. This was checked against $596{,}268{,}874$ enumerated
slots with no disagreement, and a parameter count reproduces $\alpha=3$
(\texttt{LEAD-ICP-xedni-function-field-sections}, \texttt{EV-XEDN-001}).
The scope is one family ($\deg b\le10$, $p\le31$). A dimension-count theorem
over all section shapes would be needed to close the route in general.

\section{Twist points carry no relations}

A curve and its quadratic twist share $x$-coordinates, and therefore share
summation polynomials. Could twist points supply relations for the curve?
They cannot. Frobenius acts as $+1$ on $E(\FF_q)$ and as $-1$ on the
embedded twist, so the twist part $B$ of any sum satisfies $2B=0$.

The study \art{research/twist_obstruction_20261026/} checked this
exactly on toy binary curves in the ECC2K-130 model (odd $n$, twist
$a\mapsto a+1$):
\begin{itemize}
\item 0 of 356{,}530{,}036 mixed sums are rational;
\item 0 violations on 4{,}954 decomposition targets at $m=2,3$;
\item the off-parity solution sets equal the predicted ones on 1{,}154 of
  1{,}154 targets;
\item positive controls found all 18{,}042 and 95 planted solutions.
\end{itemize}
A heuristic that ignored the obstruction predicted that 74--89\% of
solutions would contain a twist-lifting coordinate; 0\% did at $m=2$.

A follow-up asked whether a factor-base subspace can \emph{avoid} the twist
altogether (\art{research/lift_dense_factor_base_20261027/}; 11 Lean
theorems, standard axioms only).
\begin{itemize}
\item \emph{Composite degree $N=dn$, $n\ge3$: yes.} On suitable
  $\FF_{2^d}$-lines the lift condition collapses to the constant
  $\operatorname{Tr}(a)$. When $\operatorname{Tr}(a)=0$ every abscissa
  lifts, giving $2^n$ relations per Gaudry--Diem trial. This applies to
  c2pnb208w1 and c2pnb368w1, but the modelled gain is $+0.50$ bits, a
  model output, not a measurement.
\item \emph{Prime degree (ECC2K-130): this construction does not exist.}
  A prime-degree field has no proper intermediate subfield, so there are no
  $\FF_{2^d}$-lines to use. This does not rule out other twist-avoiding
  subspaces. At the small prime degrees tested, all-lift subspaces of
  dimension 4--5 exist, though fewer than for density-paired random sets
  (exact all-lift 4-subspace counts $0.921\pm0.018$ of the null at $N=11$,
  permutation $p=2\cdot10^{-5}$). This is toy evidence. Whether useful
  twist-avoiding subspaces exist at cryptographic prime degree is
  unresolved.
\end{itemize}

\section{What the obstructions leave open}

None of the obstructions above concerns the cost of \emph{solving} the
decomposition system, nor the linear algebra. Chapter~\ref{ch:falldegree}
shows that the solving cost is genuinely open. The algorithms of
Part~\ref{part:algorithms} attack the remaining free terms.
